% SPDX-License-Identifier: GPL-3.0-or-later
% BlenderMesh2STEP user guide.  Build with:  latexmk -pdf user-guide.tex
\documentclass[11pt,a4paper]{article}

\usepackage[margin=2.6cm]{geometry}
\usepackage[T1]{fontenc}
\usepackage[utf8]{inputenc}
\usepackage{lmodern}
\usepackage{microtype}
\usepackage{parskip}
\usepackage{booktabs}
\usepackage{enumitem}
\usepackage{xcolor}
\usepackage{fancyhdr}
\usepackage{listings}
\usepackage[colorlinks=true, linkcolor=blue!50!black, urlcolor=blue!50!black]{hyperref}
\usepackage[most]{tcolorbox}

\definecolor{panelgray}{gray}{0.95}
\definecolor{warnred}{RGB}{180,60,40}
\definecolor{okgreen}{RGB}{40,130,60}

\newtcolorbox{tipbox}{colback=okgreen!6, colframe=okgreen!70!black,
  title=Tip, fonttitle=\bfseries, boxrule=0.6pt, left=6pt, right=6pt}
\newtcolorbox{warnbox}{colback=warnred!6, colframe=warnred,
  title=Important, fonttitle=\bfseries, boxrule=0.6pt, left=6pt, right=6pt}

\lstset{basicstyle=\ttfamily\small, backgroundcolor=\color{panelgray},
  frame=single, framerule=0pt, breaklines=true, columns=fullflexible}

\newcommand{\ui}[1]{\textsf{\textbf{#1}}}   % a button / panel label
\newcommand{\version}{0.11.0}

\pagestyle{fancy}
\fancyhf{}
\fancyhead[L]{\small BlenderMesh2STEP User Guide}
\fancyhead[R]{\small v\version}
\fancyfoot[C]{\thepage}

\title{\Huge BlenderMesh2STEP\\[0.4em]
  \Large Turn Blender meshes into real CAD STEP solids\\[0.8em]
  \large User Guide --- version \version}
\author{}
\date{2026-07-11}

\begin{document}
\maketitle
\thispagestyle{empty}

\begin{abstract}
\noindent
BlenderMesh2STEP is a Blender 4.2+ extension that converts triangle meshes into
genuine STEP AP242 B-rep solids --- the kind FreeCAD, Fusion~360, SolidWorks
and Onshape open as editable, measurable, manufacturable geometry. You select
the faces of a feature, the extension fits the exact analytic surface (plane,
box, cylinder, cone, sphere, torus, fillet) by least squares --- typically to
machine precision on clean Blender-authored meshes --- and the exporter writes a
true solid, with real holes and pockets cut by boolean operations. This guide
covers installation, the reverse (mesh $\to$ CAD) workflow, forward parametric
modeling, exporting, and verifying the result.
\end{abstract}

\tableofcontents
\clearpage

%=====================================================================
\section{Introduction}

\subsection{The problem}

A mesh has thrown away the \emph{intent} you want back: a
20\,mm cylinder and a 64-sided prism are identical triangle soups. Exporting
STL (or tessellated STEP) hands a machine shop a faceted blob they cannot
measure, fillet, or edit.

BlenderMesh2STEP does not guess --- \textbf{you point, it reconstructs}. You
tell it ``these faces are a cylinder,'' and it recovers the exact analytic
surface by least-squares fitting. The result is a true boundary-representation
(B-rep) solid with real \texttt{CYLINDRICAL\_SURFACE} entities, not an
approximation. On clean meshes the fits land at $\sim$$10^{-8}$ RMS ---
effectively exact.

\subsection{Two workflows}

\begin{description}[style=nextline]
  \item[Reverse (mesh $\to$ CAD)] Select the faces of one feature in Edit
    Mode, fit a primitive, repeat per feature, export. This is the main
    workflow and the subject of Section~\ref{sec:reverse}.
  \item[Forward (build $\to$ CAD)] Add parametric STEP-primitive solids
    directly from the \ui{Build --- STEP Primitives} panel. These are
    STEP-exact \emph{by construction}; the viewport mesh is only a preview.
    See Section~\ref{sec:forward}.
\end{description}

Both workflows feed the same \emph{feature stack} and the same exporter, and
can be mixed freely --- e.g.\ fit a base plate from a mesh, then add a
forward-built cylinder as a \emph{Subtract} feature to drill a hole.

\subsection{Two export backends}
\label{sec:backends}

\begin{description}[style=nextline]
  \item[OCCT kernel (recommended)] Uses OpenCASCADE --- the kernel FreeCAD is
    built on --- installed on demand from the panel (Section~\ref{sec:occt}).
    This is the path that can \emph{cut holes} (booleans), merge bodies into
    one watertight solid, heal gaps, and validate the result (per-solid
    volume and topology checks).
  \item[Pure-Python writer (fallback)] Zero dependencies. Writes real analytic
    AP242 solids, but has \textbf{no boolean kernel}: subtractive features
    cannot be cut from the part (Section~\ref{sec:purepython}).
\end{description}

\begin{warnbox}
If your part has holes, pockets, counterbores, or must be a single watertight
body, install the OCCT kernel. Without it a ``hole'' can only be marked,
skipped, or written as filled material --- never actually cut.
\end{warnbox}

%=====================================================================
\section{Installation}

\subsection{Requirements}

\begin{itemize}[nosep]
  \item Blender \textbf{4.2 or newer} (the extension uses the 4.2+ extension
    format).
  \item No other dependencies for fitting and the pure-Python STEP writer.
  \item Optionally, the OCCT kernel ($\sim$100\,MB, one click) for booleans,
    watertight sewing, and validation.
\end{itemize}

\subsection{Installing the extension}

\begin{enumerate}[nosep]
  \item Download \texttt{reverse\_mesh-\version.zip} from the project's GitHub
    Releases page.
  \item In Blender: \ui{Edit} $\to$ \ui{Preferences} $\to$ \ui{Get Extensions}
    $\to$ (top-right dropdown) \ui{Install from Disk\ldots} and pick the zip.
  \item Enable it if it is not enabled automatically. The \ui{Reverse} tab
    appears in the 3D-viewport sidebar (press \texttt{N}).
\end{enumerate}

\subsection{Installing the OCCT kernel}
\label{sec:occt}

In the \ui{Reverse} tab, open \ui{Fitted Features} $\to$ \ui{Export} and click
\ui{Install OCCT}. This downloads the \texttt{cadquery-ocp} binding into the
extension's own directory (it does not touch your system Python). When it is
ready the panel shows \ui{OCCT ready}. All boolean and watertightness options
in the export dialog are marked \emph{(OCCT only)} and silently do nothing on
the pure-Python path.

%=====================================================================
\section{Quick start: the 60-second workflow}

\begin{enumerate}[nosep]
  \item Select your mesh object, enter \textbf{Edit Mode}, open the
    \ui{Reverse} tab in the sidebar (\texttt{N}).
  \item Choose a \ui{Primitive} (or leave \ui{Auto-detect}) and a \ui{Role}:
    \emph{Add} for material, \emph{Subtract} for a hole/pocket cutter.
  \item Select the faces of \emph{one} feature. Click one face of a wall and
    use \ui{Select Similar} to grow the selection across the whole surface
    (it stops at sharp edges).
  \item Click \ui{Fit Primitive to Selection}. The fit lands in the
    \ui{Fitted Features} stack with its RMS error; the optional heatmap
    colours the faces by deviation. Repeat for each feature.
  \item Click \ui{Export STEP (AP242)}, review the options
    (Section~\ref{sec:export}), export, and read the \ui{Validation Report}
    panel to confirm volumes and watertightness.
\end{enumerate}

%=====================================================================
\section{The reverse workflow in detail}
\label{sec:reverse}

\subsection{Primitives you can fit}

\begin{table}[h]
\centering
\small
\begin{tabular}{@{}lll@{}}
\toprule
\textbf{Primitive} & \textbf{Recovers} & \textbf{Notes} \\
\midrule
Plane    & flat face with extent            & exported as a bounded planar patch \\
Box      & oriented cuboid (rotation too)   & fit from several flat faces at once \\
Cylinder & axis, radius, height             & the workhorse for holes and bosses \\
Cone     & apex, half-angle, radii          & full cones and frustums \\
Sphere   & centre, radius                   & \\
Torus    & axis, major + minor radius       & best on full rings \\
Fillet   & edge round $\to$ partial cylinder & radius + arc extent; arcs $<330^\circ$ \\
Extrude  & prism: planar line/arc profile + height & select the whole prism incl.\ caps \\
Auto     & best of the above                & normal-agreement gate + simplicity tie-break \\
\bottomrule
\end{tabular}
\caption{Fittable primitives. All fits report RMS and maximum error.}
\end{table}

The \textbf{Extrude} fit recovers the archetypal Blender-authored mechanical
part: a planar outline of lines and arcs swept along an axis (brackets,
plates, slots, bosses). Select the \emph{entire} prism — side walls and both
end caps — and the fitter recovers the axis, the height, and an exact
line/arc profile (tangent arcs such as slot ends are recognised as true
arcs, not chains of facets). It exports as a genuine B-rep solid on both
backends and can act as a \emph{Subtract} cutter (a milled pocket).

\ui{Auto-detect} tries every primitive and keeps the best fit, preferring the
\emph{simplest} primitive when several explain the selection equally well
(a flat patch becomes a plane, not a giant sphere). The runner-up candidates
are shown for the active feature so you can see how close the call was.

\subsection{Selecting a feature's faces}

\begin{itemize}[nosep]
  \item \ui{Select Similar} grows the selection from the active face across
    edges whose face-normal angle is below the \emph{crease angle} --- one
    click on a cylindrical wall selects the whole wall and stops at the caps.
  \item \ui{Segment regions} (in the main panel) splits a multi-surface
    selection into smooth-connected regions and fits each separately --- e.g.\
    a whole cube selection becomes 6 planes; a full cylinder becomes side +
    2 caps.
\end{itemize}

\subsection{Fit quality controls}

\begin{description}[style=nextline]
  \item[Fit-quality heatmap] Colours the selected faces green $\to$ red by
    deviation from the fitted surface, so a stray chamfer or mis-picked face
    is obvious instead of buried in an average.
  \item[Reject outliers (RANSAC)] Trims faces that deviate from the consensus
    surface and refits on the rest. Use when a selection is slightly dirty
    (an odd triangle, the start of a chamfer).
  \item[Snap dimensions] Rounds fitted radii/lengths to a nice value
    (e.g.\ $19.98 \to 20.0$) when they are already very close, so the STEP
    carries manufacturable numbers. Presets 0.1 / 0.5 / 1.0 or a custom step.
  \item[Tolerance] The RMS threshold (as a fraction of region size) above
    which the fit is flagged as poor.
\end{description}

\subsection{Propagate Pattern: fit one hole, get the rest}

Fit \emph{one} hole (a subtractive cylinder), then click
\ui{Propagate Pattern}. The extension scans the mesh for every matching hole
--- same radius and axis direction within tolerance --- and fits them all with
the seed's settings. Bolt circles, linear arrays and mirrored holes are
recognised and labelled (circular / linear / scattered).

\subsection{Whole-mesh auto-decompose (experimental)}

Two heavier, whole-mesh paths exist for when you want a first pass without
hand-selecting anything. Both are marked red in the UI because they optimise
over the \emph{entire} mesh and can take a while:

\begin{description}[style=nextline]
  \item[Auto-Decompose Whole Mesh] Surface decomposition: sweeps
    coarse$\to$fine crease angles, builds competing primitive hypotheses, and
    selects a set by global energy minimisation (data fit + coverage +
    a per-primitive complexity cost + boundary smoothness). Faces no
    primitive explains can be kept as a \texttt{Reverse\_Leftover} faceted
    patch so the exported STEP still contains the complete part.
  \item[Auto-Decompose Solid (Boolean)] Volumetric decomposition: samples the
    interior volume and greedily covers it with inscribed spheres, cylinders
    and boxes, producing an additive union of solids. Export with the OCCT
    \ui{Merge into one solid} option.
\end{description}

Treat both as a convenience to correct, not a promise: review the resulting
feature set, delete bad members, and re-fit where needed.

%=====================================================================
\section{The feature stack}
\label{sec:stack}

Every fit and every forward-built primitive lands in \ui{Fitted Features}.
Each entry shows its boolean role (\texttt{+}/\texttt{--}), kind, a summary
with dimensions, and its RMS (or ``exact'' for built primitives). Prefixes mark
machine-generated sets: \texttt{[A]} surface auto-decompose, \texttt{[S]}
solid decompose, \texttt{[B]} forward-built.

From the stack you can, non-destructively:

\begin{itemize}[nosep]
  \item \textbf{Reorder} features (arrows). Order matters with
    \emph{Booleans in stack order} at export: an Add placed after a cut
    refills it (e.g.\ a boss inside a pocket).
  \item \textbf{Re-fit} a feature from its remembered source faces after
    editing the mesh.
  \item \textbf{Toggle Add / Subtract}, and for subtractive features choose
    \ui{Through} (cutter overshoots both ends so the hole opens cleanly) or
    \ui{Blind} (keeps the pocket depth exact; the open end is auto-detected).
  \item \textbf{Tag threads}: on a cylinder/cone, enter e.g.\ \texttt{M8x1.25}.
    The designation is written into the STEP product name and the PMI sidecar.
  \item \textbf{Hole presets}: on a subtractive cylinder choose
    \ui{Counterbore} (recess radius + depth) or \ui{Countersink} (recess
    radius + angle); the recess is cut as an extra coaxial cutter on the OCCT
    path.
  \item \textbf{Select} the feature's generated object in the viewport;
    \textbf{delete} single features or \textbf{clear} whole sets.
\end{itemize}

The stack persists in the \texttt{.blend} file; the authoritative parameters
of each feature live on its generated object, so features survive reloads and
renames of the session list.

%=====================================================================
\section{Forward modeling: Build --- STEP Primitives}
\label{sec:forward}

The \ui{Build} panel adds solids that are STEP-exportable by construction:
\textbf{Box, Cylinder, Cone (frustum), Sphere, Torus, Extrude (N-gon prism)}.
The exact analytic parameters are stored on the object; the viewport mesh is
just a preview tessellation. For the extruded prism, pick the side count when
adding; the circumradius and height stay live-editable and the profile follows.

\begin{itemize}[nosep]
  \item Pick a primitive and a \ui{Role} (Add or Subtract --- a forward-built
    cylinder makes a perfect drill), then \ui{Add Primitive}.
  \item With the object active, its dimensions appear as live-editable fields;
    editing one regenerates the mesh immediately.
  \item If you move/rotate the object, that transform is carried into the
    export. If you \emph{scale} it in Object Mode, the panel offers
    \ui{Bake Scale into Parameters} --- click it so the stored dimensions match
    what you see. Non-uniform scale on curved primitives is flagged: it is
    not representable as an analytic STEP surface.
  \item If the mesh drifts out of sync with the stored parameters (e.g.\ you
    edited vertices by hand), the panel warns and offers
    \ui{Rebuild from Parameters}.
\end{itemize}

%=====================================================================
\section{Exporting STEP}
\label{sec:export}

Click \ui{Export STEP (AP242)} in the \ui{Fitted Features} panel. The file
dialog carries the export options:

\subsection{Backend and boolean options}

\begin{table}[h]
\centering
\small
\begin{tabular}{@{}p{4.2cm}p{9.6cm}@{}}
\toprule
\textbf{Option} & \textbf{Meaning} \\
\midrule
Backend & \ui{Auto} (OCCT if installed, else pure Python), or force either. \\
Merge into one solid & Fuse all Add solids into a single body (OCCT). \\
Booleans in stack order & Apply Add/Subtract in feature-stack order so a
  later Add refills a cut. Off = fuse all Adds first, then cut (OCCT). \\
Cutter overshoot & Extend cutters by this fraction at each end so through
  holes open cleanly on coplanar faces (OCCT). \\
Auto-stitch shared edges & Unify coincident faces so abutting features share
  real edges (OCCT). \\
Make watertight + Sew tolerance & Sew all faces, build a closed solid, heal
  gaps up to the tolerance; reports remaining open edges (OCCT). \\
Cutters (pure Python) & What to do with Subtract features when there is no
  kernel: \ui{Include marked} (red \texttt{cutter:} reference solids ---
  default), \ui{Skip}, or \ui{Include as solids} (holes appear filled). \\
\bottomrule
\end{tabular}
\caption{Backend/boolean export options.}
\end{table}

\subsection{Units and scale}
\label{sec:units}

The \ui{Unit} option sets the length unit declared inside the STEP file
(mm / m / inch); it defaults to your scene's display unit. \ui{Scale from
Scene units} derives the coordinate scale from Blender's unit settings ---
e.g.\ a metric scene at Unit Scale 0.001 with STEP unit mm gives
\texttt{1 BU $\to$ 1 mm}. The dialog always shows the \emph{effective scale}
live; check it before exporting. Choose \ui{Manual} to force a factor.
Scenes with unit system \emph{None} pass coordinates through unchanged.

\subsection{Scope and extras}

\begin{itemize}[nosep]
  \item \ui{Selected only} exports just the selected Reverse objects.
  \item \ui{Feature set} filters by group: All / Manual / Built / Auto /
    Solid --- useful when an auto-decompose set coexists with hand fits.
  \item \ui{Include leftover patches} writes the unexplained-face patches so
    the STEP contains the complete part (as faceted planar faces).
  \item \ui{Write PMI sidecar} adds \texttt{.pmi.json} and \texttt{.pmi.csv}
    next to the STEP with every feature's dimensions (radii, diameters,
    lengths, angles, threads) and pairwise relationships (axis angles, hole
    spacing) --- handy for inspection sheets.
  \item \ui{Embed semantic PMI} embeds AP242 \texttt{DIMENSIONAL\_SIZE}
    entities so CAD reads diameters/lengths as queryable PMI (pure-Python
    writer only).
\end{itemize}

\subsection{The validation report}

After an OCCT export, the \ui{Validation Report} sub-panel lists every solid
with its volume and validity, plus a watertightness line (open-edge count if
not). Read it every time: \emph{valid + watertight + plausible volume} is
your evidence the file will survive a CAD import and a machine shop's quote.
The pure-Python path cannot validate; it says so instead of pretending.

\subsection{Limits of the pure-Python writer}
\label{sec:purepython}

The built-in writer emits real analytic solids (box, cylinder, cone, sphere,
torus), bounded plane patches, trimmed partial-cylinder fillet patches, and
faceted leftover patches --- with per-feature colours. It cannot: perform
booleans (no holes), merge bodies, sew to watertight, or validate. Overlapping
Add bodies are written as separate interpenetrating solids in one part.

%=====================================================================
\section{Verification and troubleshooting}

\subsection{Verifying a STEP file}

\begin{itemize}[nosep]
  \item Open it in FreeCAD: it should appear as solid bodies, not a mesh.
    \ui{Part} workbench $\to$ \ui{Check geometry} should report no errors for
    watertight exports.
  \item Measure a known dimension (a hole diameter, a plate thickness) ---
    the point of this tool is that those numbers are exact.
  \item Compare the reported volume in the \ui{Validation Report} against an
    expectation ($\pi r^2 h$ for a cylinder, etc.).
\end{itemize}

\subsection{Common issues}

\begin{table}[h]
\centering
\small
\begin{tabular}{@{}p{5.6cm}p{8.2cm}@{}}
\toprule
\textbf{Symptom} & \textbf{Likely cause / fix} \\
\midrule
Part imports 1000$\times$ too big/small &
  Unit mismatch. Check the \emph{effective scale} line in the export dialog
  (Section~\ref{sec:units}). \\
Holes are filled in the imported part &
  Exported on the pure-Python path with \ui{Include as solids}, or OCCT not
  installed. Install OCCT; keep \ui{Cutters} on \ui{Include marked}. \\
Hole doesn't break through the surface &
  Increase \ui{Cutter overshoot}, or set the feature's cut mode to
  \ui{Through}. \\
Watertight pass reports open edges &
  Leftover/fillet patches only meet fitted surfaces within the fit deviation.
  Raise \ui{Sew tolerance} slightly, or exclude leftovers. \\
Fit has a high RMS &
  The selection includes faces from another surface. Use the heatmap to spot
  them, or enable \ui{Reject outliers}. \\
Auto-detect picks the wrong primitive &
  Select more of the surface (partial patches are ambiguous), or choose the
  primitive explicitly --- the runner-up list shows how close the call was. \\
Torus/fillet fit is approximate &
  Known limitation on partial patches; select as much of the ring/round as
  possible. Fillet arcs must be $<330^\circ$. \\
STEP file is huge &
  Leftover patches carry many triangles. Remesh, or disable \ui{Include
  leftover patches}. \\
\bottomrule
\end{tabular}
\caption{Troubleshooting quick reference.}
\end{table}

\begin{tipbox}
Model with intent in Blender and the fits will be exact: keep flat things
flat, use enough segments on curved surfaces, apply scale
(\texttt{Ctrl+A}) before fitting, and prefer one feature per surface.
\end{tipbox}

%=====================================================================
\section{Roadmap}

Development continues along the plan in \texttt{PLAN.md} at the repository
root: an OCCT-first experience, extruded-profile and revolve recovery
(genuine sketch-based solids), fillets and chamfers applied as real blends,
multi-body STEP with colours, and --- longer term --- links from the feature
stack into Blender's simulation systems and an optional feature-based CNC
toolpath module.

\appendix
\section{Running the test suite}

From the repository root:

\begin{lstlisting}
# Pure-NumPy core (no Blender needed)
python3 reverse_mesh/tests/test_fitting.py
python3 reverse_mesh/tests/test_step.py
python3 reverse_mesh/tests/test_decompose.py
python3 reverse_mesh/tests/test_solidfit.py

# OCCT kernel round-trip (needs cadquery-ocp)
python3 reverse_mesh/tests/test_occ_export.py

# Blender integration (headless)
blender --background --python reverse_mesh/tests/blender_smoke.py
\end{lstlisting}

\section{License}

BlenderMesh2STEP is free software under the GPL-3.0-or-later license.

\end{document}
